\section{Using the data rows as source pivots}
\label{sec:borrowed-source-pivots}

This section proves a replacement finite bit word and describes its conditional
histogram under the inherited two-stage partial-swap transfer. It does not
independently prove the all-size machine simulation or multiplication assembly.
All scalar coefficients in this section lie in $\mathbb F_2$.

\paragraph{Finite-word identity.}
Let $X=(X_S)_{S\in\binom{[h]}3}$ be the source data bank, $Y$ the target bank,
and $Z$ the remaining arbitrary auxiliary inputs. In the positive producer,
replace each of the $v=\binom h3$ leaf-pivot roles by the corresponding data row
$X_S$. Keep all middle CNOTs and every other physical role. Write $M$ for this
invertible middle word on $(X,Z)$, and $J$ for its scatter map, including both
retained-centre and side contributions. The unchanged coefficient construction
says
\[
 JM(X,Z)=X+DZ
\]
for a fixed binary matrix $D$. Apply the following operations in chronological
order:
\[
 Y\gets Y+DZ;\qquad M;\qquad Y\gets Y+JM(X,Z);\qquad M^{-1}.
\]
In the scatter term, $J$ reads the current middle-word outputs. Thus the two
occurrences of $DZ$ cancel, $X$ and $Z$ are restored, and the exact public map is
\[
 (X,Y,Z)\longmapsto(X,Y+X,Z).
\]
No invertibility of an auxiliary-only restriction of $M$ is assumed. The
arbitrary auxiliary contribution is computed explicitly, rather than treated
as zero.

\paragraph{Forward frame chronology.}
Use the inherited factor-space labels, with base frame $D_0$, full stage frame
$D_1$, source line $\langle t_S\rangle$, retained point-star span $U_c$, and
side-target frame $t_T^\perp$. The early correction uses only $Y$ and $Z$, both
at $D_0$; it never reads $X$ at that frame. During the middle word, $X_S$ starts
at its source-line frame, which is precisely the old leaf pivot's first
incidence label. Every later incidence therefore has the same nested source
span as in the original producer. Every remaining auxiliary role is explicitly placed at
$D_0$, including roles whose column in $D$ vanishes. Thus their selected
outer edges remain exactly the inherited ones, rather than being folded into
a larger first middle jump. No middle CNOT involves $Y$.

After completing $M$, first perform all retained-centre contributions while
every $Y_T$ is still at $D_0$. Read each retained total from its existing $U_c$
role through the usual fresh temporary copied from $U_c$ to $D_0$, paying rank
$h-1$. These values have already been computed by $M$; the original $X$ rows
need not still contain their input values. Next grow each $Y_T$ to its target
frame and perform the side scatters. Finally grow all $X$ and $Z$ roles to
$D_1$ and execute $M^{-1}$ pointwise there. The retained-centre roles never
return. Performing side scatters before centre reads would require a forbidden
return of $Y$ to $D_0$ and is explicitly rejected by the checker.

\paragraph{The opposite shear with the same chronology.}
Transpose every elementary CNOT while retaining its chronological position.
Since a CNOT is its own inverse, the resulting total matrix is the
inverse transpose of the forward word. The forward public shear is itself an
involution, so the public map becomes
\[
 (X,Y,Z)\longmapsto(X+Y,Y,Z).
\]
A transposed CNOT touches exactly the same two roles at the same common frame;
therefore every established frame incidence remains valid. The transposed
retained-centre bridge first gathers its target rows into a fresh temporary at
$D_0$, moves that temporary from $D_0$ to $U_c$, and adds it to the retained
role there. Its charged rank is again $h-1$. The source-span interpretation of
the new intermediate scalar values is not needed: compatibility follows from
the literal common-frame gates and their transpose identities.

For the second tensor stage put $P=\langle t_{a_1}\rangle$ and
$B=P^\perp\subset F$. The increasing local frame $U$ is lifted to
\[
 D_U=(B\otimes F)\mathbin{\perp}(P\otimes U).
\]
The inherited stage-two entrances place $X$ at $D_{\langle t_{a_2}\rangle}$
and $Y$ at $D_0$. Insert an explicit identity vertex at this source-line frame
for every $X$ role, even when its first actual middle CNOT occurs at a larger
frame. This retains the original data-corner edge from $F\otimes
\langle t_{a_2}\rangle$ before the unchanged local source chain starts;
no entrance is merged with the first middle growth. The same-chronology inverse-transpose word then ends at
$D_F=F^{\otimes2}$ on $X$ and
$D_{t_{a_2}^\perp}=u_a^\perp$ on $Y$, exactly the required endpoints.
The first stage uses the original forward lift tensored with the fixed
second-factor source line. Thus the two scalar shears and the charged endpoint
correction are unchanged.

\paragraph{A sparse implementation of the early correction.}
The literal matrix $D$ is fixed, so the early correction is always a finite
pointwise XOR program. Alternatively, supply $v$ freshly allocated zero
streams as the old leaf pivots, run the original middle word on these and $Z$,
scatter at $D_0$, and undo that middle word. This computes $DZ$, restores $Z$,
and returns the fresh streams to zero. It uses $2|M|+|J|$ pointwise gates.
For the inverse-transpose realization, the private fresh streams may receive
$Y$-dependent values; clear them at $D_0$ after that early block. Their contents
are not arbitrary public inputs and no address change is involved. Both
orientations are checked on every public input basis vector. The direct dense
realization is also checked and avoids dependence on this storage optimization.

\paragraph{Profile accounting.}
Let $R,W,H=(n_w)$ be the original retained-total role count, volume denominator,
and child histogram at $m=h^2$, $N=v^2$. Then
\[
 R'=R-v,\qquad W'=W-2N.
\]
A borrowed source data row follows the old pivot's path after its source line,
replacing the old direct data edge of rank $h-1$. For each of the $2N$ absorbed
source pivots, remove its selected auxiliary edge with child widths
$\{m-2h,h\}$, its initial singleton entrance, and that old direct data edge.
Write $\mathcal I(h-1)=\{h-2,1\}$ for the inherited inner profile. Thus
\[
 H'=H-2N\bigl(\{m-2h,h\}+\{1\}+\mathcal I(h-1)\bigr).
\]
The removed child widths sum to $m$. Consequently
\[
 s'=s-2Nm,\qquad W'm-s'=Wm-s.
\]
Every remaining child has its original width and its original frame idempotent;
no new generic-basis condition is introduced by this change. The borrowed data
rows replace existing direct edges; their internal growth is not counted twice.
The retained-centre loss and endpoint correction remain fully charged.

\paragraph{Exact dimension-23 result and limits.}
For the paired-tree, reassociated producer,
\[
 h=23,\quad v=1771,\quad R'=38306,\quad
 W'=141952734,\quad s'=75091652097,
\]
with unchanged deficit $1344189$. The complete candidate histogram has the
exact rational exponential-moment certificate
\[
 a_b=\frac{38238149}{10^{12}}=0.000038238149,
\]
compared with $36943733/10^{12}$ for that producer before borrowing the source
rows. This is approximately a $3.5\%$ improvement in this bit saving, not a
measured runtime improvement and not the requested $2^{-10}$ witness.

The finite replay covers all $41848$ arbitrary public scalar basis inputs at
$h=23$, the forward and same-chronology transposed words, every monotone middle
source-frame history, and the centre-before-side ordering. It also covers
$h=5,6,7,8$. Omitting a nonzero early cancellation gate fails the full dirty
identity; moving side scatters before centre reads fails the frame chronology.
The dimension-23 middle word has $73047$ CNOTs, and its sparse complete
realization has $313440$ pointwise gates plus initialization and clearing of
private temporary storage. The remaining global fixed-tape, recursive,
ordered-affine, and analytic interfaces are inherited obligations, not facts
established solely by these finite tests.
